\documentclass[10pt,a4paper]{amsart}
\usepackage[margin=19mm]{geometry}
\usepackage{amsmath,amssymb,mathtools}
\usepackage[scr=rsfs,cal=euler]{mathalfa}
\usepackage{xfrac}
\usepackage[unicode,hidelinks,bookmarks=false]{hyperref}
\newcommand{\ie}{i.e.\ }
\newcommand{\cf}{cf.\ }
\newcommand{\resp}{respectively\ }
\newcommand{\Subsection}{\subsection}
\providecommand{\nobreakdash}{\nobreak\hbox{-}}
\setlength{\emergencystretch}{2em}
\multlinegap0pt
\usepackage{bm}
\usepackage{bbm}
\usepackage{dsfont}
\usepackage{xspace}
\usepackage{cancel}
\usepackage{mathtools}
\usepackage{amsthm}
\makeatletter
\@ifundefined{theorem}{\newtheorem{theorem}{Theorem}}{}
\@ifundefined{lemma}{\newtheorem{lemma}[theorem]{Lemma}}{}
\makeatother
\newcommand{\ZZ}{\mathbb{Z}}
\newcommand{\mono}{\rightarrowtail}
\newcommand{\spec}{\mathrm{Spec\hskip .5mm }}
\title[Endomorphisms of Equivariant Algebraic $K$-theory]{Endomorphisms of Equivariant Algebraic $K$-theory}
\author{K. Arun Kumar, Girja S Tripathi}
\date{Source: arXiv:2304.02544v2 (CC BY 4.0)}
\begin{document}
\maketitle

\noindent\emph{Source and licence.} Endomorphisms of Equivariant Algebraic $K$-theory. Version arXiv:2304.02544v2 (CC BY 4.0), distributed under the Creative Commons Attribution 4.0 International licence, as recorded in the arXiv licence metadata for this exact version and shown on the abstract page https://arxiv.org/abs/2304.02544v2. Full rights notice: licenses/arxiv-2304.02544-CC-BY-4.0.txt.\par\smallskip

\noindent\textit{Editorial scope.} Counted unit: the Theorem whose source label is \texttt{th:ontopi0} in the article (source0.tex lines 642--644, source label \texttt{th:ontopi0}), delivered together with the article's own complete original proof of it (source0.tex lines 645--664, one \texttt{proof} environment of 20 lines). No mathematics is written, repaired or re-derived: statement and proof are the article's own text line for line. Required context retained uncounted: Equivalent Category of Summands in Regular Representations (lines 193--196); eq:maptopi0 (lines 634--636). Every definition, display and label of those lines is retained unchanged. Omitted from this package and not redistributed: 1--192, 197--633, 637--641, 665--893. Nothing inside a retained excerpt is omitted or altered: every retained body is delivered exactly as the article's lines are written, so no omission receipt of any kind accompanies this package. Citations: the retained excerpts carry no citation command at all, so no bibliography entry of the article is transcribed and no reference is redistributed. Reference adaptation: none was needed and none was made; every reference inside the delivered unit resolves inside it, so no unresolved reference or citation is produced. Printed slips kept literal: no printed word, symbol, hypothesis or number of the counted statement or of its proof was altered. Layout and class: the article's own class is \texttt{amsart} with packages including tikz, tikz-cd, geometry, latexsym, amssymb, graphicx, hyperref, mathrsfs, todonotes; the wrapper is the lane's pinned \texttt{amsart} editorial wrapper at 10pt on A4, which restates the theorem-like environments the retained text needs and adds the article's own macro definitions that the retained text needs (\texttt{\textbackslash ZZ}, \texttt{\textbackslash mono}, \texttt{\textbackslash spec}), with extra wrapper packages (bm, bbm, dsfont, xspace, cancel, mathtools). The delivered numbering is the unit's own. Assets: no retained span contains a figure environment, a graphics-inclusion command, a PDF-page inclusion, a table or a TikZ picture; no image file is copied into this package, the package declares no asset dependency and no image file is redistributed with it. Licence: version arXiv:2304.02544v2 of this article is distributed under the Creative Commons Attribution 4.0 International licence, as recorded in the arXiv licence metadata for this exact version and shown on the abstract page https://arxiv.org/abs/2304.02544v2. A full rights notice is delivered at licenses/arxiv-2304.02544-CC-BY-4.0.txt. Characters: every retained excerpt is pure ASCII, so the delivered engine, XeLaTeX with its default Latin Modern text font, reports no missing character.
\begin{lemma}
\label{Equivalent Category of Summands in Regular Representations}
For an affine $G$-scheme $\spec R$, under our assumption in the paper that the order of the group is invertible in $R$, every $G$-equivariant finitely generated projective $R$-module is a summand in a finite direct sum of regular representations.
\end{lemma}

	\begin{equation}\label{eq:maptopi0}
	Hom_{M^G(S)}(X,\ZZ\times Gr_G)\to K_0(G,X) 
	\end{equation}

	\begin{theorem}\label{th:ontopi0}
		For all affine $X\in Sm^G_S$, the map in $\mathrm(\ref{eq:maptopi0})$ is onto.
	\end{theorem}

	\begin{proof}
     It is sufficient to show this for $G$-connected schemes. Let $X\cong \spec R$ be affine and $G$-connected. There is an epimorphism of $G$-vector bundles $s:\rho \to R$ given by $s(\sum_g r_ge_g)=\sum r_g$. Let $I$ denote the kernel of this map. The morphism $Hom_{M^G(S)}(X,\{i\}\times Gr(r\textsl{g},\rho^{\oplus n}))\to K_0(G,X)$ is given by
		\begin{equation}\label{loc:map}
		(i,M\mono \rho^{\oplus n})\mapsto [M]-r[\rho]+i[R] 
		\end{equation}  
		 for each $r,n\in\mathbb{N}$ and $i\in\ZZ$ as discussed above. Let $[P]-[Q]\in K_0(G,\spec R)$. By Lemma~\ref{Equivalent Category of Summands in Regular Representations}, $Q$ is a direct summand of $\rho^{\oplus n}$ for some $n$. Therefore $[Q]=n[\rho]-[Q']$ for some $G$-vector bundle $Q'$ and we have 
		 \[[P]-[Q]=[P]-n[\rho]+[Q']=[P\oplus Q']-n[\rho]=[P']-n[\rho]\]
		 in $K_0(G, \spec R)$. If rank of $P'$ is $p=b\textsl{g}+r$ with $r<\textsl{g}$ then we can add and subtract $(\textsl{g}-r)$ copies of $[R]$ to get 
		 \[ [P']-n[\rho]=[P'\oplus R^{\oplus \textsl{g}-r}]-n[\rho]-(\textsl{g}-r)[R] \]
		 where $P'\oplus R^{\oplus \textsl{g}-r}$ then has rank a multiple of $\textsl{g}$. Therefore every element in $K_0(G,\spec R)$ is of the form $[M]-n[\rho]-i[R]$ where rank of $M$ is a multiple of $\textsl{g}$. By \ref{loc:map} it is enough to show that any element of this form is equal to some $[M']-n'[\rho]-i'[R]$ with rank of $M=n'\textsl{g}$. Let $rank(M)=m\textsl{g}$ we then have two cases, $m<n$ and $m>n$. If $m<n$ then 
		 \begin{multline*}
		 [M]-n[\rho]-i[R]=[M]+(n-m)\textsl{g}[R]-(n-m)\textsl{g}[R]-n[\rho]-i[R]\\=[M\oplus R^{\oplus (n-m)\textsl{g}}]-n[\rho]-(n\textsl{g}-m\textsl{g}+i)[R]
		 \end{multline*}
		 which is in the desired form. When $m>n$ then 
		 \begin{multline*}
		 	[M]-n[\rho]-i[R]=[M]+(m-n)\textsl{g}[I]-(m-n)\textsl{g}[I]-n[\rho]-i[R]\\
		 	=[M\oplus I^{\oplus (m-n)\textsl{g}}]-(m-n)\textsl{g}[\rho]-n[\rho]+((m-n)\textsl{g}-i)[R]\\=[M\oplus I^{\oplus (m-n)\textsl{g}}]-((m-n)\textsl{g}+n)[\rho]+((m-n)\textsl{g}-i)[R]
		 \end{multline*}  
		 which is of the desired form.
	\end{proof}

\end{document}
